\chapter{Variance Swaps and Volatility Derivatives}\label{ch:dv:variance-swaps-and-volatility-derivatives}

A \omterm{def:dv:variance-swaps-and-volatility-derivatives:vs}{variance swap} has no strike in price. Its only strike is a volatility, squared, and it pays the difference
between that number and the variance the index actually realises. It looks like a bet that needs a model of
volatility. It needs none. Its hedge is a fixed strip of every out-of-the-money option on the index, each
weighted by one over its strike squared, plus a position in the index that is rebalanced every day. The
volatility index is that strip's price, quoted as a volatility. This chapter builds the replication, the
index formula and its errors, and the convexity that separates a swap on volatility from a swap on variance.
It then covers futures and options on the volatility index, and the jumps and truncated strikes that break
the replication.

\section{Variance swaps and their replication}

\begin{definition}[Variance swap]\label{def:dv:variance-swaps-and-volatility-derivatives:vs}
A \emph{variance swap}\index{variance swap} is a forward contract on realised variance. At expiry $T$ the long
side receives $N_{\mathrm{var}}\bigl(\sigma_R^2-K^2\bigr)$, where
$\sigma_R^2=\frac{252}{n}\sum_{i=1}^n\ln^2(S_{t_i}/S_{t_{i-1}})$ is the annualised realised variance of daily
closes, with no mean subtracted, and $K$ is the strike, quoted as a volatility.
\end{definition}

\begin{definition}[Variance notional]\label{def:dv:variance-swaps-and-volatility-derivatives:varnot}
The \emph{variance notional}\index{variance notional} $N_{\mathrm{var}}$ is the payment per unit of variance
(per ``variance point'' when variance is quoted in percent squared).
\end{definition}

\begin{definition}[Vega notional]\label{def:dv:variance-swaps-and-volatility-derivatives:veganot}
The \emph{vega notional}\index{vega notional} of a \omterm{def:dv:variance-swaps-and-volatility-derivatives:vs}{variance swap} is $N_{\mathrm{vega}}=2KN_{\mathrm{var}}$: the
approximate payment for one volatility point of realised volatility above the strike, near the strike.
Trades are quoted in \omterm{def:dv:greeks-and-the-hedging-pnl:greeks}{vega} notional and settled in \omterm{def:dv:variance-swaps-and-volatility-derivatives:varnot}{variance notional}.
\end{definition}

The payoff is convex in realised volatility. With a \omterm{def:dv:variance-swaps-and-volatility-derivatives:veganot}{vega notional} of 100\,000 per point and a strike of
22.1\%, the \omterm{def:dv:variance-swaps-and-volatility-derivatives:varnot}{variance notional} is $100\,000/(2\times22.1)=2\,262$ per variance point. A realised volatility
of 30\% pays $2\,262\times(30^2-22.1^2)=931\,000$ rather than the 790\,000 a linear payoff would give, and
15\% costs 596\,000 rather than 710\,000. The long side gains more than proportionally when volatility spikes.
That convexity is why a \omterm{def:dv:variance-swaps-and-volatility-derivatives:vs}{variance swap} is replicable and a \omterm{def:dv:variance-swaps-and-volatility-derivatives:volswap}{volatility swap} is not.

\begin{proposition}[Replication]\label{prop:dv:variance-swaps-and-volatility-derivatives:rep}
If the index moves continuously, with any volatility process, then
\[
\int_0^T\sigma_t^2\,dt=2\int_0^T\frac{dS_t}{S_t}-2\ln\frac{S_T}{S_0},
\]
and, with $F$ the forward and zero rates,
\[
-2\ln\frac{S_T}{F}=-2\,\frac{S_T-F}F+\int_0^F\frac2{K^2}(K-S_T)^+\,dK+\int_F^\infty\frac2{K^2}(S_T-K)^+\,dK.
\]
Realised variance is therefore the payoff of a static strip of out-of-the-money options with weights
$2/K^2$, a forward, and a dynamic position of $2/S_t$ in the index. The fair strike is
\[
K^2T=2\int_0^F\frac{P(K)}{K^2}\,dK+2\int_F^\infty\frac{C(K)}{K^2}\,dK.
\]
\end{proposition}
\begin{proof}
Itô's formula (One Quant Book 4, chapter 3) for $\ln S$ gives
$d\ln S_t=dS_t/S_t-\frac12\sigma_t^2\,dt$; integrate and rearrange. For the second identity, any twice
differentiable payoff satisfies
$f(S)=f(F)+f'(F)(S-F)+\int_0^Ff^{\prime\prime}(K)(K-S)^+dK+\int_F^\infty f^{\prime\prime}(K)(S-K)^+dK$ (Taylor's formula with integral
remainder), and $f(S)=-2\ln(S/F)$ has $f^{\prime\prime}(K)=2/K^2$. The dynamic leg $2\int dS/S$ has zero price, and so
has the forward, which leaves the strip.
\end{proof}

The proof never used a model for $\sigma_t$. That is the product's appeal: the strike is read from option
prices alone. The dealer who sells the swap buys the strip once, trades the index every day and is
hedged, whatever volatility does, as long as the index does not jump. The strip's weights also say what
the strike is made of (\cref{fig:dv:variance-swaps-and-volatility-derivatives:strip}). On chapter 9's
surface the one-year strike is 22.1\%, against an at-the-money volatility of 19.2\%. Puts carry 72\% of the
variance, strikes below 80 carry 24\%, and strikes above 120 carry 1.5\%. A \omterm{def:dv:variance-swaps-and-volatility-derivatives:vs}{variance swap} on a skewed index is
largely a position in its downside.

\begin{omfigure}
\begin{tikzpicture}
\begin{groupplot}[group style={group size=2 by 1,horizontal sep=1.6cm},
  width=6.6cm,height=5cm,
  tick label style={font=\small},label style={font=\small},grid=major,grid style={gray!25},table/col sep=comma]
\nextgroupplot[xlabel={strike},ylabel={share of the variance (\%)},xmin=25,xmax=165,ymin=0,ymax=17,
  title={\small one-year strip, strikes every 5}]
\addplot[omThm,ybar,bar width=3.5pt,fill=omThm!40] table[x=k,y=share]{figdata/derivatives/14-variance-swaps-and-volatility-derivatives/contrib.csv};
\draw[gray,dashed] (axis cs:100,0) -- (axis cs:100,17);
\nextgroupplot[xlabel={index at expiry},ylabel={payoff},xmin=40,xmax=170,ymin=-1.2,ymax=1.9,
  title={\small $-2\ln(S_T/F)$ and a strip},
  legend columns=1,legend cell align=left,legend style={font=\footnotesize,at={(0.5,-0.32)},anchor=north,draw=none,fill=none}]
\addplot[omDef,very thick] table[x=s,y=log]{figdata/derivatives/14-variance-swaps-and-volatility-derivatives/replication.csv};
\addplot[omThm,thick,dashed] table[x=s,y=strip]{figdata/derivatives/14-variance-swaps-and-volatility-derivatives/replication.csv};
\legend{log payoff,{strip, strikes 60 to 140 every 10}}
\end{groupplot}
\end{tikzpicture}

\omcaption{Left: each strike's share of the one-year variance-swap strike on chapter 9's surface: the puts
below the forward (dashed) carry 72\%. Right: the log payoff and its replication by nine strikes. The strip is
piecewise linear, exact at the strikes up to the spacing error, and linear beyond the last strikes, where it
stops replicating. Data: the chapter's code.}\label{fig:dv:variance-swaps-and-volatility-derivatives:strip}
\end{omfigure}

A \omterm{def:dv:variance-swaps-and-volatility-derivatives:vs}{variance swap} is marked to market by splitting it in time. After a time $t$ with realised variance
$\sigma_{R,t}^2$, the expected final variance is the time-weighted average
$\bigl(t\,\sigma_{R,t}^2+(T-t)K_{t,T}^2\bigr)/T$ of the realised part and the fair strike of the remaining
part. The value is $N_{\mathrm{var}}$ times the discounted difference from $K^2$. A one-year swap struck at
22.1\% that has realised 25\% over six months, with six-month variance-swap volatility at 20.1\%, expects
$0.5\times25^2+0.5\times20.1^2=514.5$ variance points against a strike of 488.4. On a \omterm{def:dv:variance-swaps-and-volatility-derivatives:varnot}{variance notional} of
2\,262 it is worth 59\,000.

\section{The log contract and the volatility-index formula}

\begin{definition}[Log contract]\label{def:dv:variance-swaps-and-volatility-derivatives:log}
The \emph{log contract}\index{log contract} pays $\ln(S_T/F)$ at $T$. Its price, $-\frac12\sigma_{\mathrm{VS}}^2T$
with zero rates, is the variance-swap strike in another unit. It is not listed, and is traded as its option
strip.
\end{definition}

The volatility index turns the strip into a published number. On a listed option chain the integral becomes a
sum over strikes, and the index formula is
\[
\sigma^2=\frac2T\sum_i\frac{\Delta K_i}{K_i^2}\,e^{RT}Q(K_i)-\frac1T\Bigl(\frac F{K_0}-1\Bigr)^2,
\]
with $Q(K_i)$ the mid-quote of the out-of-the-money option at $K_i$, $K_0$ the strike at or just below the
forward, and $\Delta K_i$ half the distance between neighbouring strikes. The last term corrects for using
the put and call at $K_0$ rather than at $F$. Between $K_0$ and $F$ the strip holds an in-the-money call as an
out-of-the-money put, and the term removes the difference to second order. Two expiries that bracket thirty
days are computed and interpolated to a constant thirty days.

\begin{dated}{2026-09}{The volatility index's methodology}\label{dat:dv:variance-swaps-and-volatility-derivatives:vix}
Cboe's methodology document for the index (version 6.0, last revised 26 February 2026) specifies the formula above.
It uses S\&P~500 options at mid-quote, keeps only options with a non-zero bid, brackets thirty days with two
expiries, standard or weekly, and states the formula's basis as the 1999 research note of Demeterfi, Derman,
Kamal and Zou. A filtering algorithm guards against quotes that widen suddenly. Futures and options on the
index settle on a special opening quotation that follows slightly different rules from the spot index.
\end{dated}

How good is the sum? It has two errors: the spacing between strikes, and the range the strikes cover, which in
the index is cut wherever bids vanish. On chapter 9's thirty-day smile, whose exact strip volatility is 16.30\%,
the index-style number errs as in \cref{tab:dv:variance-swaps-and-volatility-derivatives:errors}. Wide spacing
biases it up, since a coarse sum over a convex function overstates it. A short range biases it down, since the
missing wings would have added variance. A strip that stops at 95 misses 1.6 to 1.8 points. The two errors can
cancel by accident, as the strip from 90 in steps of 2.5 shows.

\begin{table}[ht]
\centering
\begin{tabular}{rrrr}
\toprule
lowest strike (highest $=200-$lowest) & $\Delta K=0.5$ & $\Delta K=1$ & $\Delta K=2.5$\\
\midrule
70 & $+0.01$ & $+0.06$ & $+0.38$\\
80 & $-0.02$ & $+0.03$ & $+0.36$\\
85 & $-0.13$ & $-0.07$ & $+0.28$\\
90 & $-0.48$ & $-0.40$ & $+0.01$\\
95 & $-1.78$ & $-1.61$ & $-0.92$\\
\bottomrule
\end{tabular}
\caption{Error of the index-style thirty-day volatility against the exact strip (16.30\%), in volatility points,
by strike range and spacing, on chapter 9's surface (spot 100). Data: the tutorial.}\label{tab:dv:variance-swaps-and-volatility-derivatives:errors}
\end{table}

\section{Volatility swaps and convexity}

\begin{definition}[Volatility swap]\label{def:dv:variance-swaps-and-volatility-derivatives:volswap}
A \emph{volatility swap}\index{volatility swap} pays $N\bigl(\sigma_R-K_{\mathrm{vol}}\bigr)$: realised
volatility, not variance, against a strike.
\end{definition}

The payoff is linear in volatility, which is what many clients want to trade, but volatility is the square
root of variance and no static strip replicates a square root of a path integral. Its fair strike is
$\E[\sigma_R]$, and Jensen's inequality puts it below the variance-swap strike $\sqrt{\E[\sigma_R^2]}$. A
second-order expansion of the square root around the mean variance $V$ gives
\[
K_{\mathrm{vol}}\approx\sqrt V-\frac{\Var(\sigma_R^2)}{8V^{3/2}},
\]
so the convexity adjustment is proportional to the variance of realised variance: it needs a model for the
\omterm{def:dv:stochastic-volatility:sv}{volatility of volatility}. It is the convexity adjustment of One Quant Book 2, chapter 8, applied to a square
root.

In Heston's model it is exact. The Laplace transform of integrated variance,
$\E[e^{-s\int_0^Tv_t\,dt}]$, has the closed form of an affine model. The identity
$\sqrt x=\frac1{2\sqrt\pi}\int_0^\infty(1-e^{-sx})s^{-3/2}ds$ turns it into $\E[\sqrt{\cdot}]$ by one
integral. The build does exactly that, and its tests check it against a simulation.

\begin{example}[Volatility or variance]\label{ex:dv:variance-swaps-and-volatility-derivatives:convexity}
Under the \omterm{def:dv:stochastic-volatility:heston}{Heston model} calibrated to chapter 9's surface in chapter 10, the one-year variance-swap strike is
21.4\% and the volatility-swap strike 19.8\%: a convexity adjustment of 1.62 volatility points. The adjustment
is 0.72 point at one month, 1.34 at three months and 1.61 at six months, then falls to 1.35 at two years and 1.10
at three (\cref{fig:dv:variance-swaps-and-volatility-derivatives:convexity}). At short expiries there is little
time for variance to vary. At long expiries mean reversion averages its variation away. The market's own
one-year strip gives 22.1\%, 0.7 point above Heston's, because the calibrated model does not reproduce the
surface's wings, which carry weight in the strip.
\end{example}

\begin{omfigure}
\begin{tikzpicture}
\begin{axis}[width=10.5cm,height=5cm,
  xlabel={expiry (years)},ylabel={strike (volatility, \%)},
  tick label style={font=\small},label style={font=\small},
  xmin=0,xmax=3.1,ymin=15,ymax=25,grid=major,grid style={gray!25},
  legend columns=3,legend cell align=left,legend style={font=\footnotesize,at={(0.5,-0.3)},anchor=north,draw=none,fill=none,/tikz/every even column/.append style={column sep=8pt}},
  table/col sep=comma]
\addplot[omDef,very thick,mark=*,mark size=1.6pt] table[x=t,y=var]{figdata/derivatives/14-variance-swaps-and-volatility-derivatives/heston_terms.csv};
\addplot[omThm,very thick,mark=square*,mark size=1.6pt] table[x=t,y=vol]{figdata/derivatives/14-variance-swaps-and-volatility-derivatives/heston_terms.csv};
\addplot[gray,thick,dashed,mark=o,mark size=1.6pt] table[x=t,y=market]{figdata/derivatives/14-variance-swaps-and-volatility-derivatives/heston_terms.csv};
\legend{variance swap (Heston),volatility swap (Heston),variance swap (market strip)}
\end{axis}
\end{tikzpicture}

\omcaption{Variance-swap and volatility-swap strikes by expiry under the \omterm{def:dv:stochastic-volatility:heston}{Heston model} calibrated to chapter 9's
surface, and the market strip's variance-swap strike. The gap between the first two is the convexity
adjustment, largest near one year. Data: the tutorial.}\label{fig:dv:variance-swaps-and-volatility-derivatives:convexity}
\end{omfigure}

Variance can be weighted and restricted, and two variants are traded. Each is replicated by the same argument
with a different function of the price.

\begin{definition}[Gamma swap]\label{def:dv:variance-swaps-and-volatility-derivatives:gamma}
A \emph{gamma swap}\index{gamma swap} pays realised variance weighted by the index level relative to its initial
value, $\frac1T\int_0^T(S_t/S_0)\,\sigma_t^2\,dt$, against a strike. It is replicated by the payoff
$\frac2{S_0}(S\ln(S/S_0)-S+S_0)$, whose second derivative is $2/(S_0K)$: its strip weights puts less than the
\omterm{def:dv:variance-swaps-and-volatility-derivatives:vs}{variance swap}'s.
\end{definition}

\begin{definition}[Corridor variance swap]\label{def:dv:variance-swaps-and-volatility-derivatives:corridor}
A \emph{corridor variance swap}\index{corridor variance swap} accrues realised variance only on days when the
index is inside a range $[L,U]$. Its strip keeps the weights $2/K^2$ on strikes in $[L,U]$ only.
\end{definition}

The \omterm{def:dv:variance-swaps-and-volatility-derivatives:gamma}{gamma swap}'s variance does not explode when the index falls, since each day's variance is multiplied by
$S_t/S_0$, which is small after a crash. That makes it a much less dangerous product to be short than the \omterm{def:dv:variance-swaps-and-volatility-derivatives:vs}{variance swap}. On
chapter 9's surface its one-year strike is 20.8\%, between the at-the-money 19.2\% and the \omterm{def:dv:variance-swaps-and-volatility-derivatives:vs}{variance swap}'s
22.1\%. A down-corridor ($U=S_0$) isolates the variance paid in falling markets, the part of the \omterm{def:dv:variance-swaps-and-volatility-derivatives:vs}{variance swap}
that the skew makes expensive.

\section{Futures and options on the volatility index}

Futures on the volatility index have traded since March 2004 and options since February 2006. Chapter 12 showed
that the index is, idealised, the square root of the average \omterm{def:dv:rough-volatility-and-forward-variance-models:fwdvar}{forward variance} over thirty days, and that its
future is $\E[\mathrm{VIX}_T]$. The future therefore lies below the forward variance-swap volatility
$\sqrt{\E[\mathrm{VIX}_T^2]}$ by a convexity gap. The gap is the same Jensen effect as the \omterm{def:dv:variance-swaps-and-volatility-derivatives:volswap}{volatility swap}'s,
applied to thirty days of variance seen from $T$.

\begin{example}[The futures curve under Heston]\label{ex:dv:variance-swaps-and-volatility-derivatives:vixcurve}
Under the calibrated \omterm{def:dv:stochastic-volatility:heston}{Heston model} the index is 17.1 today. The three-month future is 18.1 against a forward
thirty-day variance-swap volatility of 20.4, the six-month 19.3 against 22.2, and the one-year 20.5 against 23.7
(\cref{fig:dv:variance-swaps-and-volatility-derivatives:vix}, left). The model computes the futures exactly:
$\mathrm{VIX}_T^2=\bar v+(v_T-\bar v)\frac{1-e^{-\kappa\Delta}}{\kappa\Delta}$ is affine in $v_T$, whose Laplace
transform is known, so the same square-root identity applies. A desk that reads the forward volatility off the
futures curve without the gap overstates the futures' fair level by one to three points.
\end{example}

\begin{definition}[VIX option]\label{def:dv:variance-swaps-and-volatility-derivatives:vixopt}
A \emph{VIX option}\index{VIX option} is a European option on the volatility index at its expiry. It is
cash-settled, and its natural underlying is the VIX future of the same expiry, so it is quoted with Black's
formula on the future.
\end{definition}

A \omterm{def:dv:variance-swaps-and-volatility-derivatives:vixopt}{VIX option} is an option on the forward-variance curve, so its smile measures the distribution of future
volatility. The market's VIX smile slopes upward: calls on high volatility are expensive. Baldeaux and Badran
point to this observed upward skew. Under the calibrated \omterm{def:dv:stochastic-volatility:heston}{Heston model} the three-month smile slopes the other way,
from 110\% at 80\% of the future to 94\% at 175\%
(\cref{fig:dv:variance-swaps-and-volatility-derivatives:vix}, right). The square-root process has a thin right
tail. Rough Bergomi's VIX smile is nearly flat (chapter 12). Fitting the VIX smile's slope with the same model as
the index smile is the joint calibration problem of chapter 12.

\begin{omfigure}
\begin{tikzpicture}
\begin{groupplot}[group style={group size=2 by 1,horizontal sep=1.6cm},
  width=6.6cm,height=5cm,
  tick label style={font=\small},label style={font=\small},grid=major,grid style={gray!25},table/col sep=comma]
\nextgroupplot[xlabel={futures expiry (months)},ylabel={volatility points},xmin=0,xmax=12,ymin=16,ymax=25,
  title={\small index futures, Heston},
  legend columns=1,legend cell align=left,legend style={font=\footnotesize,at={(0.5,-0.32)},anchor=north,draw=none,fill=none}]
\addplot[omThm,very thick,mark=*,mark size=1.4pt] table[x=months,y=future]{figdata/derivatives/14-variance-swaps-and-volatility-derivatives/vix_curve.csv};
\addplot[omDef,thick,dashed] table[x=months,y=fwd]{figdata/derivatives/14-variance-swaps-and-volatility-derivatives/vix_curve.csv};
\legend{future,forward 30-day variance-swap volatility}
\nextgroupplot[xlabel={strike (\% of the future)},ylabel={implied volatility (\%)},xmin=75,xmax=180,ymin=85,ymax=120,
  title={\small three-month VIX options, Heston}]
\addplot[omThm,very thick,mark=*,mark size=1.6pt] table[x=strike,y=vol]{figdata/derivatives/14-variance-swaps-and-volatility-derivatives/vix_smile.csv};
\end{groupplot}
\end{tikzpicture}

\omcaption{The volatility index under the calibrated \omterm{def:dv:stochastic-volatility:heston}{Heston model}. Left: futures lie below the forward
variance-swap volatility by a convexity gap that grows with expiry. Right: the three-month smile of options on
the index slopes down, the opposite of the market's upward slope. Data: the tutorial.}\label{fig:dv:variance-swaps-and-volatility-derivatives:vix}
\end{omfigure}

\begin{definition}[Variance risk premium]\label{def:dv:variance-swaps-and-volatility-derivatives:vrp}
The \emph{variance risk premium}\index{variance risk premium} is the difference between the variance-swap strike
and the expected realised variance under the real-world measure, $K^2-\E^{\mathbb P}[\sigma_R^2]$ (some authors
use the opposite sign). It is what the seller of variance earns on average for bearing the risk of variance
spikes.
\end{definition}

Where the premium is positive, a seller of \omterm{def:dv:variance-swaps-and-volatility-derivatives:vs}{variance swaps} earns it on average and pays it back, with interest,
in crashes, when realised variance exceeds the strike many times over. Its size, sign and variation across
underlyings are empirical questions. Carr and Wu (2009) is the classic study of them with synthetic \omterm{def:dv:variance-swaps-and-volatility-derivatives:vs}{variance
swaps}, and One Quant Book 6 measures them. The distinction matters for this chapter: the replication argument
gives the strike under the pricing measure, and says nothing about whether selling it pays.

\section{Skew, jumps and the limits of replication}

The replication holds if the index moves continuously. Over a day with log-return $r$, the hedge (the strip
plus the daily position $2/S$) earns $2(e^r-1-r)$ while the swap's leg pays $r^2$. The difference is
\[
2(e^r-1-r)-r^2=\frac{r^3}3+\frac{r^4}{12}+\cdots,
\]
negligible for ordinary days and not for crashes. A fall of 20\% in log costs the hedged seller
$0.04-0.0375=0.0025$ of variance per unit of notional on that one day. In a model with jumps the fair strike is
therefore not the strip: $\E[\sigma_R^2]T=K_{\mathrm{strip}}^2T-\E\bigl[\sum(2(e^r-1-r)-r^2)\bigr]$. With negative
jumps the correction is positive, and the strip underprices the \omterm{def:dv:variance-swaps-and-volatility-derivatives:vs}{variance swap}.

\begin{example}[The strip under Merton's jumps]\label{ex:dv:variance-swaps-and-volatility-derivatives:jumps}
With chapter 13's Merton model (3.16 jumps a year of mean log-size $-5.8\%$ and dispersion 4.6\%), each jump
leaves the hedge short by $-0.000182$ of variance on average ($\E[J^3]/3=-0.000188$ to leading order). The fair
strike exceeds the strip's by 0.00058 of variance a year, or 0.18 volatility point. On 100\,000 three-month paths,
a seller who strikes at the strip and hedges daily loses 0.17 point on average (0.18 up to simulation error), with a standard deviation of 0.40
point and a one-in-a-hundred loss of 1.9 points. In a diffusion of the same variance the same hedge has a
standard deviation of 0.014 point (\cref{fig:dv:variance-swaps-and-volatility-derivatives:hedge}).
\end{example}

\begin{omfigure}
\begin{tikzpicture}
\begin{axis}[width=10.5cm,height=5cm,ymode=log,
  xlabel={loss of a short 3-month variance swap, hedged (volatility points)},ylabel={probability of a larger loss},
  tick label style={font=\small},label style={font=\small},
  xmin=0,xmax=3,ymin=1e-5,ymax=1,grid=major,grid style={gray!25},
  legend columns=2,legend style={font=\footnotesize,at={(0.5,-0.3)},anchor=north,draw=none,fill=none,/tikz/every even column/.append style={column sep=8pt}},
  table/col sep=comma,unbounded coords=jump]
\addplot[omDef,very thick,dashed] table[x=loss,y=diffusion]{figdata/derivatives/14-variance-swaps-and-volatility-derivatives/hedge_tail.csv};
\addplot[omThm,very thick] table[x=loss,y=jumps]{figdata/derivatives/14-variance-swaps-and-volatility-derivatives/hedge_tail.csv};
\legend{diffusion,Merton jumps}
\end{axis}
\end{tikzpicture}

\omcaption{The replication's residual. A short three-month \omterm{def:dv:variance-swaps-and-volatility-derivatives:vs}{variance swap} struck at the strip and hedged with
the strip and a daily position of $2/S$ in the index, on 100\,000 paths: the probability of losing more than a
given number of volatility points. In a diffusion no path loses more than a tenth of a point; with jumps
each crash day leaves a loss, and one path in a hundred loses 1.9 points. Data: the tutorial.}\label{fig:dv:variance-swaps-and-volatility-derivatives:hedge}
\end{omfigure}

Three other limits matter in practice.
\begin{itemize}
\item \textbf{The strikes run out.} The strip exists only where options are quoted. In a crash the index falls
through the lowest strike, beyond which the hedge is linear while the payoff keeps growing. The truncation error
in \cref{tab:dv:variance-swaps-and-volatility-derivatives:errors} then turns from a pricing bias into a loss. Martin records that the large price jumps of
2008 and 2009 disrupted volatility-derivatives markets and caused the single-name variance-swap market to dry up completely.
\item \textbf{Caps.} A \omterm{def:dv:variance-swaps-and-volatility-derivatives:vs}{variance swap} can carry a cap on realised variance, set at a multiple of the strike. The
cap is a short call on variance, priced with a model of \omterm{def:dv:stochastic-volatility:sv}{volatility of volatility}, not with the strip.
\item \textbf{Discrete sampling and dividends.} Daily closes, the absence of a mean in the estimator, and the
dividend-adjusted forward all move the fair strike by small, computable amounts that the build's
mark-to-market must use consistently.
\end{itemize}

\section{Tutorial: replicate a variance swap}\label{tut:dv:variance-swaps-and-volatility-derivatives:tut}

\begin{tutorial}
\textbf{Goal.} Price a \omterm{def:dv:variance-swaps-and-volatility-derivatives:vs}{variance swap} from a strip on a synthetic surface, compute the index-style number and its
errors, value \omterm{def:dv:variance-swaps-and-volatility-derivatives:volswap}{volatility swaps} and index futures exactly under Heston, and simulate the hedge with and without
jumps. \textbf{End state:} the four figures, the table and the numbers of the weekend problem.
\begin{tutsteps}
\item \textbf{The index formula}, with $K_0$, the averaged option at $K_0$ and the forward correction:
\omcode{code/firm/varswap/firm_varswap.py}{20}{32}{The volatility-index formula on one expiry.}\label{lst:dv:variance-swaps-and-volatility-derivatives:index}
\item \textbf{The jump error} of the log-contract hedge, one line that decides the hedge's residual:
\omcode{code/firm/varswap/firm_varswap.py}{74}{78}{Shortfall of the hedge on a return $r$.}\label{lst:dv:variance-swaps-and-volatility-derivatives:jump}
\item \textbf{Heston's \omterm{def:dv:variance-swaps-and-volatility-derivatives:volswap}{volatility swap}} by the Laplace transform of integrated variance and the square-root
identity:
\omcode{code/firm/varswap/firm_varswap.py}{92}{114}{Exact volatility-swap strike under Heston.}\label{lst:dv:variance-swaps-and-volatility-derivatives:volswap}
\item \textbf{Run} \texttt{dv\_varswap.index\_errors()}, \texttt{heston\_terms()}, \texttt{vix\_curve()},
\texttt{hedge\_pnl()} and \texttt{fig\_varswap.py}.
\end{tutsteps}
\textbf{What to change next.} Replace the square-root identity by the second-order approximation and measure its
error at one year; cut the strip where the model's put price falls below 0.05 (a ``zero bid'') and recompute the
thirty-day number; price a down-corridor \omterm{def:dv:variance-swaps-and-volatility-derivatives:vs}{variance swap} with $U=100$.
\end{tutorial}

\section{Build: variance and volatility swaps}\label{bld:dv:variance-swaps-and-volatility-derivatives:varswap}

\begin{build}
\bfield{Purpose} The miniature firm's pricer and mark-to-market for \omterm{def:dv:variance-swaps-and-volatility-derivatives:vs}{variance swaps}, and the volatility-swap and
index-future values that its volatility desk quotes against.
\bfield{Interface} \texttt{index\_variance(strikes, calls, puts, fwd, t, rate)}; \texttt{strip\_from\_vols};
\texttt{log\_payoff\_strip}; \texttt{realised\_variance(prices)}; \texttt{variance\_notional(vega\_notional,
strike\_vol)}; \texttt{mark\_to\_market(var\_notional, strike\_vol, realised\_var, t\_elapsed,
fair\_var\_remaining, t\_remaining, df)}; \texttt{jump\_error(r)}; \texttt{vol\_swap\_approx};
\texttt{heston\_mean\_variance}, \texttt{heston\_vol\_swap}, \texttt{heston\_vix\_future}.
\bfield{Rules} Strikes are quoted in volatility and settled in variance; the realised-variance estimator is the
contract's (daily closes, 252 days, no mean); the strip uses the forward, not the spot; every model value is
checked against a simulation in the tests.
\bfield{Acceptance tests} \texttt{code/firm/varswap/tests/}: a flat smile's index number is its volatility,
with the forward between strikes; the strip reproduces the log payoff inside its range; the conventions
(notional conversion, mark-to-market); the jump error behaves like $r^3/3$; Heston's Laplace transforms equal one
at zero, the \omterm{def:dv:variance-swaps-and-volatility-derivatives:volswap}{volatility swap} tends to the \omterm{def:dv:variance-swaps-and-volatility-derivatives:vs}{variance swap} as the \omterm{def:dv:stochastic-volatility:sv}{volatility of volatility} vanishes and matches a
simulation, and index futures lie below the forward.
\bfield{Stretch} Capped \omterm{def:dv:variance-swaps-and-volatility-derivatives:vs}{variance swaps} under Heston; corridor and \omterm{def:dv:variance-swaps-and-volatility-derivatives:gamma}{gamma swaps} from a listed chain; the special
opening quotation's rules for settlement.
\end{build}

\begin{omsources}
\item K.~Demeterfi, E.~Derman, M.~Kamal and J.~Zou, ``More than you ever wanted to know about volatility swaps'',
Goldman Sachs Quantitative Strategies Research Notes (March 1999).
\item Cboe Global Markets, ``Cboe Volatility Index Methodology'', version 6.0 (2026).
\item P.~Carr and R.~Lee, ``Volatility derivatives'', \emph{Annual Review of Financial Economics} 1 (2009) 319--339.
\item P.~Carr and L.~Wu, ``Variance risk premiums'', \emph{Review of Financial Studies} 22(3) (2009) 1311--1341.
\item I.~Martin, ``Simple variance swaps'', NBER Working Paper 16884 (2011).
\item J.~Baldeaux and A.~Badran, ``Consistent modeling of VIX and equity derivatives using a 3/2 plus jumps
model'', arXiv 1203.5903 (2012).
\end{omsources}

\section{Exercises}

\begin{exercise}[$\star$]\label{exo:dv:variance-swaps-and-volatility-derivatives:1}
A client wants 50\,000 of \omterm{def:dv:variance-swaps-and-volatility-derivatives:veganot}{vega notional} on a \omterm{def:dv:variance-swaps-and-volatility-derivatives:vs}{variance swap} struck at 20\%. What is the \omterm{def:dv:variance-swaps-and-volatility-derivatives:varnot}{variance notional}, and what
does the swap pay if realised volatility is 25\%?
\end{exercise}

\begin{exercise}[$\star$]\label{exo:dv:variance-swaps-and-volatility-derivatives:2}
Why does a \omterm{def:dv:variance-swaps-and-volatility-derivatives:vs}{variance swap} need a daily trade in the index as well as the option strip?
\end{exercise}

\begin{exercise}[$\star$]\label{exo:dv:variance-swaps-and-volatility-derivatives:3}
Show that the \omterm{def:dv:variance-swaps-and-volatility-derivatives:gamma}{gamma swap}'s replicating payoff $\frac2{S_0}(S\ln(S/S_0)-S+S_0)$ has second derivative $2/(S_0S)$,
and explain why its strike is below the \omterm{def:dv:variance-swaps-and-volatility-derivatives:vs}{variance swap}'s on a downside-skewed index.
\end{exercise}

\begin{exercise}[$\star\star$]\label{exo:dv:variance-swaps-and-volatility-derivatives:4}
After nine months a one-year \omterm{def:dv:variance-swaps-and-volatility-derivatives:vs}{variance swap} struck at 22\% has realised 18\%, and the three-month variance-swap
volatility is 16\%. Value it per unit of \omterm{def:dv:variance-swaps-and-volatility-derivatives:varnot}{variance notional} (variance in percent squared, zero rates).
\end{exercise}

\begin{exercise}[$\star\star$]\label{exo:dv:variance-swaps-and-volatility-derivatives:5}
Use the second-order approximation to estimate the volatility-swap strike when the variance-swap strike is 20\%
and the standard deviation of realised variance is 0.02.
\end{exercise}

\begin{exercise}[$\star\star$]\label{exo:dv:variance-swaps-and-volatility-derivatives:6}
Why does the convexity adjustment in \cref{ex:dv:variance-swaps-and-volatility-derivatives:convexity} peak near one
year?
\end{exercise}

\begin{exercise}[$\star\star\star$]\label{exo:dv:variance-swaps-and-volatility-derivatives:7}
\emph{Coding.} Compute the thirty-day index-style number with strikes every 1 from 80 to 120 and then drop the
strikes whose out-of-the-money option is worth less than 0.05. How much does the number fall?
\end{exercise}

\begin{exercise}[$\star\star\star$]\label{exo:dv:variance-swaps-and-volatility-derivatives:8}
\emph{Find the flaw.} ``We sell \omterm{def:dv:variance-swaps-and-volatility-derivatives:vs}{variance swaps} at the strip's price and hold the strip and the daily delta; the
book is fully hedged, so the premium we charge over the strip is pure profit.''
\end{exercise}

\section{Problem: Vol or Variance}

\begin{problem}[{Weekend problem --- the price of a square root}]\label{pb:dv:variance-swaps-and-volatility-derivatives:1}
A client asks for one-year quotes on a \omterm{def:dv:variance-swaps-and-volatility-derivatives:vs}{variance swap} and on a \omterm{def:dv:variance-swaps-and-volatility-derivatives:volswap}{volatility swap} on the index of chapter 9's surface.
The desk has the surface, the Heston calibration of chapter 10, and the build.

\medskip\noindent\textbf{Part I --- The \omterm{def:dv:variance-swaps-and-volatility-derivatives:vs}{variance swap}.}
\begin{enumerate}
\item Give the one-year variance-swap strike from the strip and the at-the-money volatility.
\item Which strikes carry the strike? Give the puts' share.
\item For a \omterm{def:dv:variance-swaps-and-volatility-derivatives:veganot}{vega notional} of 100\,000, give the \omterm{def:dv:variance-swaps-and-volatility-derivatives:varnot}{variance notional} and the payoff at 30\% realised volatility.
\item What would the thirty-day index-style number lose with a strip that stops at 95?
\item Give the one-year gamma-swap strike.
\end{enumerate}

\medskip\noindent\textbf{Part II --- The \omterm{def:dv:variance-swaps-and-volatility-derivatives:volswap}{volatility swap} under Heston.}
\begin{enumerate}[resume]
\item Give Heston's one-year variance-swap strike, and explain the gap to the strip's.
\item Give the one-year volatility-swap strike.
\item Give the convexity adjustment at one month, one year and three years.
\item Which Heston parameter drives the adjustment, and in which direction?
\item Why is the \omterm{def:dv:variance-swaps-and-volatility-derivatives:volswap}{volatility swap} not replicable with a static strip?
\end{enumerate}

\medskip\noindent\textbf{Part III --- The index and its options.}
\begin{enumerate}[resume]
\item Give the index level and the three-month and one-year futures under Heston.
\item Give the convexity gap at three months.
\item Describe the three-month VIX smile under Heston.
\item How would the desk hedge a \omterm{def:dv:variance-swaps-and-volatility-derivatives:volswap}{volatility swap}?
\item What premium over the strip would you charge a client who buys the \omterm{def:dv:variance-swaps-and-volatility-derivatives:vs}{variance swap}, if the index can jump?
\end{enumerate}

\medskip\noindent\textbf{Part IV --- Judgement.}
\begin{enumerate}[resume]
\item The client says the \omterm{def:dv:variance-swaps-and-volatility-derivatives:volswap}{volatility swap} ``should be priced at the variance-swap strike, since both are
volatility''. Answer in two sentences.
\item Which model assumption does the convexity adjustment depend on most?
\item What reserve would you hold on a short \omterm{def:dv:variance-swaps-and-volatility-derivatives:volswap}{volatility swap} hedged with \omterm{def:dv:variance-swaps-and-volatility-derivatives:vs}{variance swaps}?
\item State the \emph{named result}: the convexity adjustment between the one-year variance-swap and
volatility-swap strikes on the chapter's surface under Heston dynamics.
\item In one sentence: why is variance, not volatility, the traded quantity?
\end{enumerate}
\end{problem}

\section{Interview questions}

\begin{interviewq}[$\star$ \iqroles{trader, researcher}]\label{iq:dv:variance-swaps-and-volatility-derivatives:1}
How do you replicate a \omterm{def:dv:variance-swaps-and-volatility-derivatives:vs}{variance swap}?
\end{interviewq}

\begin{interviewq}[$\star$ \iqroles{trader}]\label{iq:dv:variance-swaps-and-volatility-derivatives:2}
Why is the variance-swap strike above the at-the-money volatility on an index?
\end{interviewq}

\begin{interviewq}[$\star\star$ \iqroles{researcher}]\label{iq:dv:variance-swaps-and-volatility-derivatives:3}
Is a \omterm{def:dv:variance-swaps-and-volatility-derivatives:volswap}{volatility swap}'s fair strike above or below the \omterm{def:dv:variance-swaps-and-volatility-derivatives:vs}{variance swap}'s? By how much, roughly?
\end{interviewq}

\begin{interviewq}[$\star\star$ \iqroles{developer}]\label{iq:dv:variance-swaps-and-volatility-derivatives:4}
Implement the volatility-index formula for one expiry. What edge cases matter?
\end{interviewq}

\begin{interviewq}[$\star\star$ \iqroles{risk, trader}]\label{iq:dv:variance-swaps-and-volatility-derivatives:5}
What happens to a dealer short a \omterm{def:dv:variance-swaps-and-volatility-derivatives:vs}{variance swap}, hedged with the strip, when the index gaps down 20\%?
\end{interviewq}

\begin{interviewq}[$\star\star\star$ \iqroles{trader, risk}]\label{iq:dv:variance-swaps-and-volatility-derivatives:6}
Why must a VIX future be priced below the forward variance-swap volatility, and how would you trade the gap?
\end{interviewq}
